% Purpose: Supplementary material for the Pipeline Evaluation Framework
% chapter presenting
% formal mathematical derivations of bipartite matching optimization,
% catalog partition
% conservation laws, continuous point-process detection theory, the
% mathematical
% ill-posedness of ROC/specificity in continuous seismic streams,
% heavy-tailed residual distribution
% modeling, and kinematic Cramér-Rao lower bounds on phase arrival uncertainty.
% Status: Complete; publication-grade reference text.
% Source-of-truth: doc/UNITS/main/Chapter4_Pipeline_Evaluation_Framework.tex,
% OGS/src/ogsmatching.py, and OGS/src/ogsevaluation.py.

\section*{Pipeline Evaluation Framework: Bipartite Matching
Algorithms, Partition Conservation, and Continuous Detection Theory}
\addcontentsline{toc}{section}{Pipeline Evaluation Framework:
  Bipartite Matching Algorithms, Partition Conservation, and
Continuous Detection Theory}
\label{sec:supp_pipeline_evaluation}

This supplementary document provides complete mathematical derivations,
algorithmic proofs, and statistical formulations supporting
\cref{chap:pipeline_evaluation}.

% =============================================================================
% 1. BIPARTITE MATCHING OPTIMIZATION AND ALGORITHMIC FORMULATION
% =============================================================================

\subsection*{Bipartite Matching Optimization for Seismic Catalogs}
\label{subsec:supp_bipartite_matching}

Automated validation of seismological pick arrivals and earthquake
event catalogs
against an authoritative analyst bulletin requires solving a constrained
maximum-cardinality minimum-weight bipartite matching problem.

\subsubsection*{Graph Formulation}
Let $\mathcal{B} \coloneqq \{b_1, \dots, b_N\}$ denote the set of reference
bulletin entities (ground-truth picks or events) and let $\mathcal{T} \coloneqq
\{\tau_1, \dots, \tau_M\}$ represent the set of candidate detections generated
by the automated computational pipeline. We construct an
edge-weighted bipartite
graph:
\begin{equation}
  \mathcal{G} \coloneqq (\mathcal{V}_{\mathcal{B}} \sqcup
  \mathcal{V}_{\mathcal{T}}, \mathcal{E}),
  \label{eq:supp_bipartite_graph}
\end{equation}
where $\mathcal{V}_{\mathcal{B}}$ and $\mathcal{V}_{\mathcal{T}}$ are disjoint
vertex sets corresponding to $\mathcal{B}$ and $\mathcal{T}$, and
$\mathcal{E} \subseteq \mathcal{V}_{\mathcal{B}} \times
\mathcal{V}_{\mathcal{T}}$
is the set of admissible matching edges.

\paragraph{Pick-Level Admissibility:}
For phase arrivals, a bulletin pick $b \in \mathcal{B}$ and automated pick
$\tau \in \mathcal{T}$ form an admissible edge $(b, \tau) \in
\mathcal{E}_{\mathrm{picks}}$
if and only if they share identical seismic stations, identical
phase classifications,
and temporal separation strictly within the acceptance gate $\omega_t$:
\begin{equation}
  \mathcal{E}_{\mathrm{picks}} \coloneqq \left\{(b, \tau) \in
    \mathcal{B} \times \mathcal{T}
    \;\middle|\;
    \operatorname{sta}(b) = \operatorname{sta}(\tau) \;\land\;
    \operatorname{pha}(b) = \operatorname{pha}(\tau) \;\land\;
  |t(\tau) - t(b)| \le \omega_t \right\},
  \label{eq:supp_pick_admissibility}
\end{equation}
where $\omega_t \coloneqq \qty{0.50}{\second}$ is the canonical
phase tolerance window.

\paragraph{Event-Level Admissibility:}
For earthquake origin solutions, candidate pair $(b, \tau) \in
\mathcal{B} \times \mathcal{T}$
forms an admissible edge $(b, \tau) \in
\mathcal{E}_{\mathrm{events}}$ if and only if
their spatio-temporal coordinates satisfy the three-dimensional
acceptance envelope:
\begin{equation}
  \mathcal{E}_{\mathrm{events}} \coloneqq \left\{(b, \tau) \in
    \mathcal{B} \times \mathcal{T}
    \;\middle|\;
    |t_0(\tau) - t_0(b)| \le \omega_{t_0} \;\land\;
    \delta h(b, \tau) \le \omega_h \;\land\;
  |z(\tau) - z(b)| \le \omega_z \right\},
  \label{eq:supp_event_admissibility}
\end{equation}
where $\delta h(b, \tau)$ represents the horizontal great-circle
epicentral distance:
\begin{equation}
  \delta h(b, \tau) \coloneqq 2 R_{\mathrm{E}} \arcsin \left(
    \sqrt{\sin^2\left(\frac{\varphi_{\tau} - \varphi_b}{2}\right) +
      \cos\varphi_b \cos\varphi_{\tau}
    \sin^2\left(\frac{\lambda_{\tau} - \lambda_b}{2}\right)}
  \right),
  \label{eq:supp_haversine}
\end{equation}
with Earth radius $R_{\mathrm{E}} \coloneqq
\qty{6371.0}{\kilo\metre}$, origin time
gate $\omega_{t_0} \coloneqq \qty{10.0}{\second}$, epicentral distance gate
$\omega_h \coloneqq \qty{25.0}{\kilo\metre}$, and focal depth gate
$\omega_z \coloneqq \qty{20.0}{\kilo\metre}$.

\subsubsection*{Linear Sum Assignment Optimization}
To resolve optimal one-to-one pairings without greedy arbitration ambiguities,
we assign each admissible edge $(b_i, \tau_j) \in \mathcal{E}$ a dimensionless
scalar cost $C_{ij} \in [0, 1]$:
\begin{equation}
  C_{ij} \coloneqq w_t \left(\frac{|t_0(\tau_j) -
  t_0(b_i)|}{\omega_{t_0}}\right)
  + w_h \left(\frac{\delta h(b_i, \tau_j)}{\omega_h}\right)
  + w_z \left(\frac{|z(\tau_j) - z(b_i)|}{\omega_z}\right),
  \label{eq:supp_assignment_cost}
\end{equation}
with normalized convex weights $w_t + w_h + w_z \equiv 1$ (for pick matching,
$w_t \equiv 1, w_h = w_z = 0$). For pairs outside the admissibility domain
($(b_i, \tau_j) \notin \mathcal{E}$), the cost is set to an infinite barrier
$C_{ij} \coloneqq \infty$.

The optimal matching matrix $\mathbf{X} \in \{0, 1\}^{N \times M}$
is the solution
to the integer linear program:
\begin{align}
  \min_{\mathbf{X}} \quad & \sum_{i=1}^N \sum_{j=1}^M C_{ij} X_{ij}
  \label{eq:supp_ilp_obj} \\
  \text{subject to} \quad & \sum_{j=1}^M X_{ij} \le 1, \quad \forall
  i \in \{1, \dots, N\},
  && \text{Each bulletin entity matched at most once}
  \label{eq:supp_ilp_row} \\
  & \sum_{i=1}^N X_{ij} \le 1, \quad \forall j \in \{1, \dots, M\},
  && \text{Each target candidate matched at most once}
  \label{eq:supp_ilp_col} \\
  & X_{ij} \in \{0, 1\}, \quad \forall (i, j).
  && \text{Integral assignment indicators} \label{eq:supp_ilp_integral}
\end{align}
Because the constraint matrix is totally unimodular, the linear
programming relaxation
($X_{ij} \ge 0$) possesses strictly integral extreme points. The
problem is solved
in polynomial time $\mathcal{O}(\max(N, M)^3)$ via the Kuhn-Munkres (Hungarian)
primal-dual algorithm.

\subsubsection*{Step-by-Step Worked Matching Example}
To illustrate the mechanics of primal-dual assignment and phase-swap isolation,
consider a synthetic cluster of three bulletin picks and four
automated candidates
recorded across stations \texttt{A1} and \texttt{A2}:
\begin{align*}
  \mathcal{B} &\coloneqq \left\{
    b_1 = (\texttt{A1}, \mathrm{P}, \qty{10.10}{\second}),\;
    b_2 = (\texttt{A1}, \mathrm{S}, \qty{14.20}{\second}),\;
    b_3 = (\texttt{A2}, \mathrm{P}, \qty{12.50}{\second})
  \right\}, \\
  \mathcal{T} &\coloneqq \left\{
    \tau_1 = (\texttt{A1}, \mathrm{P}, \qty{10.20}{\second}),\;
    \tau_2 = (\texttt{A1}, \mathrm{S}, \qty{14.40}{\second}),\;
    \tau_3 = (\texttt{A1}, \mathrm{P}, \qty{14.30}{\second}),\;
    \tau_4 = (\texttt{A3}, \mathrm{P}, \qty{10.15}{\second})
  \right\}.
\end{align*}
Let the temporal tolerance be $\omega_t \coloneqq
\qty{0.50}{\second}$, with station \texttt{A3}
excluded from the valid station inventory
$\mathcal{S}_{\text{valid}} = \{\texttt{A1}, \texttt{A2}\}$.

\begin{enumerate}
  \item \textbf{Domain Pre-Filtering}: Candidate $\tau_4$ belongs to
    station \texttt{A3}
    $\notin \mathcal{S}_{\text{valid}}$. It is partitioned directly
    into the skipped set:
    \begin{equation}
      \tau_4 \in \mathrm{PicksSP}.
    \end{equation}
    Candidate $\tau_4$ is excised before matrix construction and
    cannot enter matching.
  \item \textbf{Admissibility Evaluation}:
    \begin{itemize}
      \item $(b_1, \tau_1)$: Identical station \texttt{A1}, phase P,
        $|\qty{10.20}{\second} - \qty{10.10}{\second}| =
        \qty{0.10}{\second} \le \omega_t \implies C_{11} =
        \frac{0.10}{0.50} = \num{0.20}$.
      \item $(b_2, \tau_2)$: Identical station \texttt{A1}, phase S,
        $|\qty{14.40}{\second} - \qty{14.20}{\second}| =
        \qty{0.20}{\second} \le \omega_t \implies C_{22} =
        \frac{0.20}{0.50} = \num{0.40}$.
      \item $(b_2, \tau_3)$: Station matches \texttt{A1}, arrival
        $|\qty{14.30}{\second} - \qty{14.20}{\second}| =
        \qty{0.10}{\second} \le \omega_t$, but phases differ
        ($\operatorname{pha}(b_2) = \mathrm{S} \ne
        \operatorname{pha}(\tau_3) = \mathrm{P}$). In standard
        phase-locked matching, $C_{23} = \infty$. In cross-phase
        swap audit, this pair is tracked as an admissible swap
        candidate with swap penalty $\epsilon_{\mathrm{SP}}$.
      \item $b_3$: No candidate exists on station \texttt{A2}
        $\implies$ Row 3 has no admissible edges ($C_{3j} = \infty,
        \forall j$).
    \end{itemize}
  \item \textbf{Cost Matrix Construction}:
    \begin{equation}
      \mathbf{C} =
      \begin{bmatrix}
        \num{0.20} & \infty & \infty \\
        \infty & \num{0.40} & \infty \\
        \infty & \infty & \infty
      \end{bmatrix}.
    \end{equation}
  \item \textbf{Optimal Solution}:
    \begin{align}
      X_{11}^* &= 1 \implies (b_1, \tau_1) \in
      \mathcal{MH}_{\mathrm{P}} \quad (\text{Matched P Pick}), \\
      X_{22}^* &= 1 \implies (b_2, \tau_2) \in
      \mathcal{MH}_{\mathrm{S}} \quad (\text{Matched S Pick}), \\
      b_3 &\implies b_3 \in \mathcal{MS}_{\mathrm{P}} \quad
      (\text{Missed Bulletin Pick}), \\
      \tau_3 &\implies \tau_3 \in \mathcal{PS}_{\mathrm{P}} \quad
      (\text{Proposed Pick Candidate}).
    \end{align}
  \item \textbf{Resulting Partition Cardinalities}:
    $\mathrm{MH} = \num{2}$, $\mathrm{MS} = \num{1}$, $\mathrm{PS} =
    \num{1}$, $\mathrm{PicksSP} = \num{1}$.
\end{enumerate}

% =============================================================================
% 2. CATALOG PARTITION CONSERVATION THEOREM AND PROOF
% =============================================================================

\subsection*{Catalog Partition Conservation and Domain Accounting Proofs}
\label{subsec:supp_partition_conservation}

A common defect in automated catalog evaluation is metric leakage caused by
inconsistent handling of spatial, temporal, and sensor boundaries.
Here we provide
formal proofs of partition disjointness and conservation.

\subsubsection*{Mathematical Partition Definitions}
Let the total set of reference bulletin entities be $\mathcal{B}$
and the total set
of target model detections be $\mathcal{T}$. Let $\Omega \subset
\mathbb{R}^3$ denote the
admissible geographic volume, $[T_{\min}, T_{\max}] \subset
\mathbb{R}$ the temporal review window,
and $\mathcal{S}_{\mathrm{valid}}$ the validated station inventory.

\begin{definition}[Boundary Operator $\partial_{\mathrm{domain}}$]
  The domain boundary predicate $\Phi_{\mathrm{domain}}(x)$
  evaluates to true if and only if
  entity $x$ lies strictly within the spatial, temporal, and station
  inventory envelope:
  \begin{equation}
    \Phi_{\mathrm{domain}}(x) \coloneqq
    \left(\mathbf{x}(x) \in \Omega\right) \;\land\;
    \left(t(x) \in [T_{\min}, T_{\max}]\right) \;\land\;
    \left(\operatorname{sta}(x) \in \mathcal{S}_{\mathrm{valid}}\right).
    \label{eq:supp_domain_predicate}
  \end{equation}
\end{definition}

\begin{definition}[Eligible Sets and Excluded Partitions]
  The bulletin and target sets are decomposed into eligible subsets
  and excluded boundary subsets:
  \begin{align}
    \mathcal{B}_{\mathrm{el}} &\coloneqq \left\{ b \in \mathcal{B}
    \;\middle|\; \Phi_{\mathrm{domain}}(b) \equiv \mathrm{True} \right\},
    && \mathcal{SM} \coloneqq \mathcal{B} \setminus
    \mathcal{B}_{\mathrm{el}} \quad (\text{Skimmed Reference Entities}),
    \label{eq:supp_def_skimmed} \\
    \mathcal{T}_{\mathrm{el}} &\coloneqq \left\{ \tau \in
      \mathcal{T} \;\middle|\; \Phi_{\mathrm{domain}}(\tau) \equiv
    \mathrm{True} \right\},
    && \mathcal{SP} \coloneqq \mathcal{T} \setminus
    \mathcal{T}_{\mathrm{el}} \quad (\text{Skipped Target Entities}).
    \label{eq:supp_def_skipped}
  \end{align}
\end{definition}

\begin{theorem}[Catalog Partition Conservation]
  \label{thm:supp_partition_conservation}
  Let bipartite matching on the eligible bipartite subgraph
  $\mathcal{G}_{\mathrm{el}}
  \coloneqq (\mathcal{B}_{\mathrm{el}} \sqcup
  \mathcal{T}_{\mathrm{el}}, \mathcal{E}_{\mathrm{el}})$
  produce matched pairs $\mathcal{MH} \subseteq
  \mathcal{B}_{\mathrm{el}} \times \mathcal{T}_{\mathrm{el}}$
  under one-to-one constraints \cref{eq:supp_ilp_row,eq:supp_ilp_col}.
  Then $\mathcal{B}$ and $\mathcal{T}$ admit unique, pairwise
  disjoint decompositions:
  \begin{align}
    \mathcal{B} &\equiv \pi_{\mathcal{B}}(\mathcal{MH}) \sqcup
    \mathcal{MS} \sqcup \mathcal{SM},
    \label{eq:supp_decomp_b} \\
    \mathcal{T} &\equiv \pi_{\mathcal{T}}(\mathcal{MH}) \sqcup
    \mathcal{PS} \sqcup \mathcal{SW} \sqcup \mathcal{SP},
    \label{eq:supp_decomp_t}
  \end{align}
  where $\pi_{\mathcal{B}}$ and $\pi_{\mathcal{T}}$ denote
  coordinate projections onto
  $\mathcal{B}$ and $\mathcal{T}$, and the cardinalities strictly
  satisfy the conservation laws:
  \begin{align}
    |\mathcal{B}| &\equiv |\mathcal{MH}| + |\mathcal{MS}| + |\mathcal{SM}|,
    \label{eq:supp_card_b} \\
    |\mathcal{T}| &\equiv |\mathcal{MH}| + |\mathcal{PS}| +
    |\mathcal{SW}| + |\mathcal{SP}|.
    \label{eq:supp_card_t}
  \end{align}
\end{theorem}

\begin{proof}
  By construction in \cref{eq:supp_def_skimmed,eq:supp_def_skipped},
  $\mathcal{B} = \mathcal{B}_{\mathrm{el}} \sqcup \mathcal{SM}$ with
  $\mathcal{B}_{\mathrm{el}} \cap \mathcal{SM} = \emptyset$.
  Under the one-to-one matching constraint \cref{eq:supp_ilp_row},
  each bulletin element
  $b \in \mathcal{B}_{\mathrm{el}}$ is either matched to exactly one
  target candidate
  $\tau \in \mathcal{T}_{\mathrm{el}}$ (in which case $b \in
  \pi_{\mathcal{B}}(\mathcal{MH})$)
  or has no assigned partner. The set of unmatched eligible bulletin
  elements is defined as:
  \begin{equation}
    \mathcal{MS} \coloneqq \mathcal{B}_{\mathrm{el}} \setminus
    \pi_{\mathcal{B}}(\mathcal{MH}).
  \end{equation}
  Because $\pi_{\mathcal{B}}(\mathcal{MH}) \subseteq
  \mathcal{B}_{\mathrm{el}}$, it follows immediately that:
  \begin{equation}
    \mathcal{B}_{\mathrm{el}} = \pi_{\mathcal{B}}(\mathcal{MH})
    \sqcup \mathcal{MS},
    \quad \text{with} \quad \pi_{\mathcal{B}}(\mathcal{MH}) \cap
    \mathcal{MS} = \emptyset.
  \end{equation}
  Substituting this into the union for $\mathcal{B}$:
  \begin{equation}
    \mathcal{B} = \left(\pi_{\mathcal{B}}(\mathcal{MH}) \sqcup
    \mathcal{MS}\right) \sqcup \mathcal{SM}
    \equiv \pi_{\mathcal{B}}(\mathcal{MH}) \sqcup \mathcal{MS}
    \sqcup \mathcal{SM}.
  \end{equation}
  Since all components are mutually disjoint, the cardinalities add directly:
  \begin{equation}
    |\mathcal{B}| = |\pi_{\mathcal{B}}(\mathcal{MH})| +
    |\mathcal{MS}| + |\mathcal{SM}|.
  \end{equation}
  Because the matching is a bijection between
  $\pi_{\mathcal{B}}(\mathcal{MH})$ and
  $\pi_{\mathcal{T}}(\mathcal{MH})$,
  $|\pi_{\mathcal{B}}(\mathcal{MH})| = |\mathcal{MH}|$,
  proving \cref{eq:supp_card_b}.

  Similarly, for the target catalog, $\mathcal{T} =
  \mathcal{T}_{\mathrm{el}} \sqcup \mathcal{SP}$
  with $\mathcal{T}_{\mathrm{el}} \cap \mathcal{SP} = \emptyset$.
  Eligible candidate detections $\tau \in \mathcal{T}_{\mathrm{el}}$
  partition into:
  (i) matched entities $\pi_{\mathcal{T}}(\mathcal{MH})$,
  (ii) cross-phase swaps $\mathcal{SW} \subseteq
  \mathcal{T}_{\mathrm{el}} \setminus \pi_{\mathcal{T}}(\mathcal{MH})$
  where an arrival of phase $P$ aligns with a bulletin arrival of
  phase $S$ (or vice versa), and
  (iii) unmatched proposed candidates $\mathcal{PS} \coloneqq
  \mathcal{T}_{\mathrm{el}} \setminus
  (\pi_{\mathcal{T}}(\mathcal{MH}) \sqcup \mathcal{SW})$.
  All four partitions are mutually disjoint by definition, yielding:
  \begin{equation}
    |\mathcal{T}| \equiv |\mathcal{MH}| + |\mathcal{PS}| +
    |\mathcal{SW}| + |\mathcal{SP}|,
  \end{equation}
  which establishes \cref{eq:supp_card_t}.
\end{proof}

\begin{corollary}[Denominator Insulation of Excluded Sets]
  \label{cor:supp_insulation}
  Neither the skimmed reference partition $\mathcal{SM}$ nor the
  skipped candidate
  partition $\mathcal{SP}$ may enter the denominator of Recall or Precision:
  \begin{equation}
    \operatorname{Recall} \coloneqq
    \frac{|\mathcal{MH}|}{|\mathcal{B}_{\mathrm{el}}| + |\mathcal{SW}|}
    = \frac{|\mathcal{MH}|}{|\mathcal{MH}| + |\mathcal{MS}| + |\mathcal{SW}|},
    \label{eq:supp_correct_recall}
  \end{equation}
  \begin{equation}
    \operatorname{Precision} \coloneqq
    \frac{|\mathcal{MH}|}{|\mathcal{T}_{\mathrm{el}}|}
    = \frac{|\mathcal{MH}|}{|\mathcal{MH}| + |\mathcal{PS}| + |\mathcal{SW}|}.
    \label{eq:supp_correct_precision}
  \end{equation}
\end{corollary}

\begin{proof}
  Admitting $\mathcal{SM}$ into the recall denominator would
  penalize the pipeline for failing
  to detect events or picks that human analysts cataloged using
  stations outside the project's
  permitted streaming inventory or outside the regional boundary.
  Conversely, admitting $\mathcal{SP}$ into the precision
  denominator would penalize the pipeline
  for producing valid picks on regional boundary stations
  intentionally excluded from downstream
  associative evaluation. Therefore, $\mathcal{SM}$ and
  $\mathcal{SP}$ must be reported strictly
  as external boundary metadata.
\end{proof}

% =============================================================================
% 3. ILL-POSEDNESS OF ROC AND CONTINUOUS POINT-PROCESS DETECTION THEORY
% =============================================================================

\subsection*{Continuous Point Process Detection and Degeneracy of ROC Analysis}
\label{subsec:supp_roc_degeneracy}

A pervasive methodological error in machine learning applications to
continuous geophysical
time series is the uncritical reporting of \ac{ROC} curves and
\ac{ROC-AUC} scores.
Here we present the mathematical foundation showing why True
Negatives and Specificity
are ill-posed and degenerate for continuous seismic streams.

\subsubsection*{Point Process Formalism}
Let continuous three-component seismic velocity waveforms
$\mathbf{v}(t) \equiv \dot{\mathbf{u}}(t) \in \mathbb{R}^3$
be acquired over continuous monitoring interval $\mathcal{I}
\coloneqq [0, T_{\mathrm{total}}]$,
with observation duration $T_{\mathrm{total}} \gg \qty{1}{\hour}$.
The true ground-truth arrival catalog is a discrete marked temporal
point process:
\begin{equation}
  \mathcal{N}_{\mathcal{B}}(t) \coloneqq
  \sum_{i=1}^{N_{\mathcal{B}}} \delta(t - t_i),
  \label{eq:supp_bulletin_process}
\end{equation}
where $t_1 < t_2 < \dots < t_{N_{\mathcal{B}}}$ denote true arrival onsets.
An automated picker generates a predicted point process:
\begin{equation}
  \mathcal{N}_{\mathcal{T}}(t) \coloneqq
  \sum_{j=1}^{N_{\mathcal{T}}} \delta(t - \tau_j).
  \label{eq:supp_target_process}
\end{equation}

\subsubsection*{The Measure-Theoretic Absence of Discrete True Negatives}
Let each arrival $t_i$ possess a physical acceptance neighborhood of
half-width $\omega_t$:
$\mathcal{U}_i \coloneqq [t_i - \omega_t, t_i + \omega_t]$.
The total positive event support in the time domain is:
\begin{equation}
  \mathcal{S}_{+} \coloneqq \bigcup_{i=1}^{N_{\mathcal{B}}} [t_i -
  \omega_t, t_i + \omega_t] \subset [0, T_{\mathrm{total}}].
  \label{eq:supp_positive_support}
\end{equation}
Under standard regional seismicity rates, events are rare phenomena.
The total Lebesgue
measure of the arrival support is bounded by:
\begin{equation}
  \mu(\mathcal{S}_{+}) \le 2 N_{\mathcal{B}} \omega_t.
  \label{eq:supp_lebesgue_pos}
\end{equation}
For example, in a \qty{30}{\day} continuous deployment
($T_{\mathrm{total}} \approx \qty{2.59e6}{\second}$)
recording $N_{\mathcal{B}} = \num{1000}$ picks with tolerance
$\omega_t = \qty{0.50}{\second}$:
\begin{equation}
  \mu(\mathcal{S}_{+}) \le \num{1000} \times \qty{1.0}{\second} =
  \qty{1.0e3}{\second},
  \quad \implies \quad
  \frac{\mu(\mathcal{S}_{+})}{T_{\mathrm{total}}} \approx
  \num{3.86e-4} \equiv \qty{0.0386}{\percent}.
\end{equation}
The complement set $\mathcal{S}_0 \coloneqq [0, T_{\mathrm{total}}]
\setminus \mathcal{S}_{+}$
represents the quiescent noise background, whose Lebesgue measure constitutes
$\qty{99.9614}{\percent}$ of the continuous record:
\begin{equation}
  \mu(\mathcal{S}_0) \equiv T_{\mathrm{total}} -
  \mu(\mathcal{S}_{+}) \approx \qty{2.589e6}{\second}.
  \label{eq:supp_lebesgue_zero}
\end{equation}

\begin{theorem}[Degeneracy of Discrete True Negatives and ROC False
  Positive Rate]
  \label{thm:supp_roc_degeneracy}
  Let the continuous time interval $[0, T_{\mathrm{total}}]$ be
  discretized into uniform
  sampling windows of duration $\Delta t > 0$. As the temporal
  resolution approaches the
  continuous limit $\Delta t \to 0$, the discrete count of True
  Negatives $\mathrm{TN}(\Delta t)$
  diverges to infinity, causing the False Positive Rate
  $\mathrm{FPR}(\Delta t)$ to vanish
  identically to zero for any finite count of false alarms:
  \begin{align}
    \lim_{\Delta t \to 0^+} \mathrm{TN}(\Delta t) &= \infty,
    \label{eq:supp_tn_divergence} \\
    \lim_{\Delta t \to 0^+} \mathrm{FPR}(\Delta t) &\equiv
    \lim_{\Delta t \to 0^+} \frac{\mathrm{FP}}{\mathrm{FP} +
    \mathrm{TN}(\Delta t)} = 0,
    \quad \forall \mathrm{FP} < \infty.
    \label{eq:supp_fpr_zero}
  \end{align}
\end{theorem}

\begin{proof}
  Let the continuous interval $[0, T_{\mathrm{total}}]$ be partitioned into
  $K(\Delta t) \coloneqq \lfloor T_{\mathrm{total}} / \Delta t
  \rfloor$ non-overlapping intervals.
  The number of windows overlapping the arrival support
  $\mathcal{S}_{+}$ is bounded above by
  $K_{+}(\Delta t) \le N_{\mathcal{B}} \lceil 2\omega_t / \Delta t \rceil$.
  The number of non-event windows in $\mathcal{S}_0$ without
  automated triggers is:
  \begin{equation}
    \mathrm{TN}(\Delta t) \ge K(\Delta t) - K_{+}(\Delta t) - \mathrm{FP}.
  \end{equation}
  Taking the limit as the sampling interval $\Delta t \to 0^+$:
  \begin{equation}
    \lim_{\Delta t \to 0^+} \mathrm{TN}(\Delta t) \ge
    \lim_{\Delta t \to 0^+} \left( \frac{T_{\mathrm{total}} - 2
    N_{\mathcal{B}} \omega_t}{\Delta t} - \mathrm{FP} \right)
    = \lim_{\Delta t \to 0^+} \left(
    \frac{\mu(\mathcal{S}_0)}{\Delta t} \right) = \infty.
  \end{equation}
  Because the number of false proposals $\mathrm{FP} \equiv
  |\mathcal{PS}| < \infty$ is finite,
  the denominator of the False Positive Rate diverges:
  \begin{equation}
    \lim_{\Delta t \to 0^+} \mathrm{FPR}(\Delta t)
    = \lim_{\Delta t \to 0^+} \frac{\mathrm{FP}}{\mathrm{FP} +
    \mathrm{TN}(\Delta t)}
    \le \lim_{\Delta t \to 0^+}
    \frac{\mathrm{FP}}{\mu(\mathcal{S}_0) / \Delta t}
    = 0.
  \end{equation}
  Consequently, Specificity $\mathrm{TNR}(\Delta t) \coloneqq 1 -
  \mathrm{FPR}(\Delta t) \to 1$
  trivially for all detector configurations.
\end{proof}

\begin{corollary}[Collapse of ROC Metric Space]
  Because $\mathrm{FPR} \equiv 0$ in the continuous limit, the
  classical \ac{ROC} curve
  collapses onto the vertical ordinate axis $(\mathrm{FPR} = 0,
  \mathrm{TPR} \in [0, 1])$.
  The area under the curve $\mathrm{ROC\text{-}AUC} \to 1.00$
  unconditionally, masking even
  catastrophic false-alarm rates. Hence, \ac{ROC} analysis is
  mathematically invalid for
  evaluating continuous seismic point-process detection.
\end{corollary}

\subsubsection*{Mathematical Superiority of Precision-Recall Space}
In contrast to \ac{ROC}, Precision-Recall (\ac{PR}) space depends
exclusively on
$\mathcal{MH}$, $\mathcal{MS}$, and $\mathcal{PS}$:
\begin{equation}
  \operatorname{Precision} \coloneqq
  \frac{|\mathcal{MH}|}{|\mathcal{MH}| + |\mathcal{PS}| + |\mathcal{SW}|},
  \quad
  \operatorname{Recall} \coloneqq
  \frac{|\mathcal{MH}|}{|\mathcal{MH}| + |\mathcal{MS}| + |\mathcal{SW}|}.
  \label{eq:supp_pr_definition}
\end{equation}
Because neither formula contains the unobservable null set
$\mathcal{S}_0$ or the divergent
$\mathrm{TN}$ count, \ac{PR} trajectories and the Precision-Recall
Area Under Curve
($\operatorname{PR-AUC} \coloneqq \int_0^1
\operatorname{Precision}(r) \,\mathrm{d}r$)
remain strictly invariant under temporal discretization and
rigorously penalize false alarms.

% =============================================================================
% 4. STATISTICAL MODELING OF HEAVY-TAILED RESIDUAL DISTRIBUTIONS
% =============================================================================

\subsection*{Statistical Modeling of Heavy-Tailed Residual Distributions}
\label{subsec:supp_residual_modeling}

Empirical seismic arrival residuals $\Delta t_k \coloneqq t(T_k) -
t(B_k)$ exhibit
significant leptokurtosis (heavy tails) caused by emergent waveform onsets,
attenuation-induced phase dispersion, and local scattering.

\subsubsection*{Theoretical Density Formulations}
To determine whether timing jitter follows classical Gaussian
dispersion or heavy-tailed
exponential behavior, we compare three candidate probability density functions:

\begin{align}
  f_G(\Delta t; \mu, \sigma)
  &\coloneqq \frac{1}{\sqrt{2\pi}\sigma} \exp\left(-\frac{(\Delta t
  - \mu)^2}{2\sigma^2}\right),
  && \text{Gaussian (Normal) Distribution} \label{eq:supp_pdf_gaussian} \\
  f_L(\Delta t; \mu, b)
  &\coloneqq \frac{1}{2b} \exp\left(-\frac{|\Delta t - \mu|}{b}\right),
  && \text{Laplace (Double-Exponential) Distribution}
  \label{eq:supp_pdf_laplace} \\
  f_C(\Delta t; \mu, \gamma)
  &\coloneqq \frac{1}{\pi \gamma \left[1 + \left(\frac{\Delta t -
  \mu}{\gamma}\right)^2\right]},
  && \text{Cauchy (Lorentzian) Distribution} \label{eq:supp_pdf_cauchy}
\end{align}
where $\mu \in \mathbb{R}$ represents the location parameter
(systematic timing bias),
$\sigma > 0$ is the standard deviation, $b > 0$ is the Laplace scale
parameter ($b = \frac{1}{\sqrt{2}}\sigma_L$),
and $\gamma > 0$ is the Cauchy half-width at half-maximum.

\subsubsection*{Breakdown of Classical Moments and Robust Scale Inversion}
Sample mean $\bar{x} \coloneqq \frac{1}{N} \sum x_i$ and sample variance
$s^2 \coloneqq \frac{1}{N-1} \sum (x_i - \bar{x})^2$ possess an
asymptotic breakdown
point of $\epsilon^* \equiv \num{0}\%$: a single erroneous pick
outlier can bias the metric
arbitrarily.

To guarantee robust estimation under emergent phase jitter, we compute the
Median Absolute Deviation (\ac{MAD}):
\begin{align}
  \operatorname{median}(\Delta t)
  &\coloneqq \Delta t_{(\lfloor (N+1)/2 \rfloor)},
  \label{eq:supp_median} \\
  \operatorname{MAD}(\Delta t)
  &\coloneqq \operatorname{median}\left(\left| \Delta t_k -
  \operatorname{median}(\Delta t) \right|\right).
  \label{eq:supp_mad}
\end{align}
Under the asymptotic normality assumption, the consistent standard
deviation estimator is:
\begin{equation}
  \hat{\sigma}_{\mathrm{MAD}} \coloneqq \num{1.4826} \times
  \operatorname{MAD}(\Delta t).
  \label{eq:supp_sigma_mad}
\end{equation}
The correction factor is derived from the standard normal cumulative
distribution $\Phi(z)$:
\begin{equation}
  \Phi^{-1}(\num{0.75}) \approx \num{0.6744897} \implies
  \frac{1}{\Phi^{-1}(\num{0.75})} \approx \num{1.4826022185}.
  \label{eq:supp_mad_factor_derivation}
\end{equation}
Similarly, the Interquartile Range (\ac{IQR}) provides an
alternative robust scale:
\begin{equation}
  \operatorname{IQR}(\Delta t) \coloneqq Q_3(\Delta t) - Q_1(\Delta t),
  \quad \implies \quad
  \hat{\sigma}_{\mathrm{IQR}} \coloneqq
  \frac{\operatorname{IQR}(\Delta t)}{\Phi^{-1}(\num{0.75}) -
  \Phi^{-1}(\num{0.25})}
  = \frac{\operatorname{IQR}(\Delta t)}{\num{1.34898}}.
  \label{eq:supp_sigma_iqr}
\end{equation}
Both $\hat{\sigma}_{\mathrm{MAD}}$ and $\hat{\sigma}_{\mathrm{IQR}}$
achieve a breakdown
point of $\epsilon^* \equiv \num{50}\%$, insulating pipeline
evaluation against severe
multi-second phase misassociations.

% =============================================================================
% 5. KINEMATIC CRAMÉR-RAO LOWER BOUND ON PICK ARRIVAL ACCURACY
% =============================================================================

\subsection*{Kinematic Cramér-Rao Lower Bound on Pick Arrival Uncertainty}
\label{subsec:supp_cramer_rao}

The ultimate theoretical limit of timing precision attainable by any
unbiased arrival
time estimator $\hat{t}_0$ is dictated by the wave kinematics and
the signal-to-noise ratio.

\subsubsection*{Kinematic Wavefield Formulation}
Let the recorded seismic particle displacement be
$\mathbf{u}(\mathbf{x}, t)$, with kinematic
particle velocity $\mathbf{v}(t) \equiv \dot{\mathbf{u}}(t) =
\frac{\mathrm{d}\mathbf{u}}{\mathrm{d}t}$
and particle acceleration $\ddot{\mathbf{u}}(t) =
\frac{\mathrm{d}^2\mathbf{u}}{\mathrm{d}t^2}$.
A single-station velocity seismogram recording an arrival at true
onset time $t_0$ is modeled as:
\begin{equation}
  v_{\mathrm{obs}}(t) \coloneqq v(t - t_0) + n(t), \quad t \in [0,
  T_{\mathrm{win}}],
  \label{eq:supp_waveform_model}
\end{equation}
where $v(t)$ is the deterministic kinematic velocity signal and
$n(t)$ is zero-mean
additive Gaussian white noise with autocovariance
$\mathbb{E}[n(t)n(\tau)] = \frac{N_0}{2}\delta(t - \tau)$.

\subsubsection*{Fisher Information Matrix and Arrival Variance Bound}
The log-likelihood functional for the arrival parameter $t_0$ given
observation $v_{\mathrm{obs}}(t)$ is:
\begin{equation}
  \ln \Lambda(v_{\mathrm{obs}} \mid t_0) \coloneqq -\frac{1}{N_0}
  \int_0^{T_{\mathrm{win}}} \left[ v_{\mathrm{obs}}(t) - v(t - t_0)
  \right]^2 \,\mathrm{d}t.
  \label{eq:supp_log_likelihood}
\end{equation}
The Fisher information scalar $J(t_0)$ is the negative expectation
of the second derivative:
\begin{align}
  J(t_0)
  &\coloneqq -\mathbb{E}\left[ \frac{\partial^2 \ln
  \Lambda}{\partial t_0^2} \right]
  \label{eq:supp_fisher_def} \\
  &= \frac{2}{N_0} \int_0^{T_{\mathrm{win}}} \left[ \frac{\partial
  v(t - t_0)}{\partial t_0} \right]^2 \,\mathrm{d}t
  \label{eq:supp_fisher_deriv} \\
  &= \frac{2}{N_0} \int_0^{T_{\mathrm{win}}} \left[ -\dot{v}(t -
  t_0) \right]^2 \,\mathrm{d}t
  \label{eq:supp_fisher_dot_v} \\
  &\equiv \frac{2}{N_0} \int_0^{T_{\mathrm{win}}} \left[ \ddot{u}(t
  - t_0) \right]^2 \,\mathrm{d}t.
  \label{eq:supp_fisher_kinematic_accel}
\end{align}

\begin{theorem}[Kinematic Acceleration Bound on Arrival Timing Uncertainty]
  \label{thm:supp_cramer_rao}
  The variance of any unbiased onset estimator $\hat{t}_0$ is
  strictly lower-bounded
  by the inverse Fisher information, which is inversely proportional
  to the total kinematic
  acceleration energy of the arriving phase:
  \begin{equation}
    \operatorname{Var}(\hat{t}_0) \ge \frac{1}{J(t_0)}
    = \frac{N_0}{2 \displaystyle\int_0^{T_{\mathrm{win}}} \left[
    \ddot{u}(t) \right]^2 \,\mathrm{d}t}
    \equiv \frac{\sigma_n^2}{\displaystyle\int_0^{T_{\mathrm{win}}}
    \left[ \dot{v}(t) \right]^2 \,\mathrm{d}t}.
    \label{eq:supp_cramer_rao_bound}
  \end{equation}
\end{theorem}

\begin{proof}
  By the Cramér-Rao inequality, for any unbiased estimator
  $\mathbb{E}[\hat{t}_0] = t_0$:
  \begin{equation}
    \operatorname{Var}(\hat{t}_0) \ge [J(t_0)]^{-1}.
  \end{equation}
  Substituting the identity $\dot{v}(t) \equiv \ddot{u}(t)$ from
  kinematic calculus into
  \cref{eq:supp_fisher_kinematic_accel} yields \cref{eq:supp_cramer_rao_bound}.
\end{proof}

\begin{remark}[Geophysical Interpretation of Timing Jitter]
  \cref{eq:supp_cramer_rao_bound} demonstrates why impulsive P
  arrivals (characterized by
    high-frequency onset and large particle acceleration
  $\ddot{u}(t)$) exhibit narrow
  residual distributions ($\sigma_{\Delta t} \le
  \qty{0.05}{\second}$), whereas emergent
  S arrivals and attenuated regional phases exhibit smaller
  acceleration integrals and
  wider timing jitter ($\sigma_{\Delta t} \ge \qty{0.20}{\second}$).
\end{remark}

% =============================================================================
% 6. SPATIAL HYPOCENTER ERROR ELLIPSOIDS AND UNCERTAINTY PROPAGATION
% =============================================================================

\subsection*{Spatial Hypocenter Error Ellipsoids and Uncertainty Propagation}
\label{subsec:supp_spatial_ellipsoid}

Spatial residuals $\Delta h_k$ and $\Delta z_k$ evaluate automated
hypocenters against
reference analyst locations.

\subsubsection*{Covariance Decomposition}
Let $\mathbf{C}_{\mathbf{x}} \in \mathbb{R}^{3 \times 3}$ denote the
spatial covariance
matrix derived from the inversion Hessian or Bayesian posterior distribution:
\begin{equation}
  \mathbf{C}_{\mathbf{x}} \coloneqq
  \begin{bmatrix}
    \sigma_{xx}^2 & \sigma_{xy}^2 & \sigma_{xz}^2 \\
    \sigma_{yx}^2 & \sigma_{yy}^2 & \sigma_{yz}^2 \\
    \sigma_{zx}^2 & \sigma_{zy}^2 & \sigma_{zz}^2
  \end{bmatrix}.
  \label{eq:supp_cov_matrix}
\end{equation}
Applying spectral decomposition to the symmetric positive-definite
covariance matrix:
\begin{equation}
  \mathbf{C}_{\mathbf{x}} = \mathbf{V} \boldsymbol{\Lambda} \mathbf{V}^\top,
  \label{eq:supp_cov_eigen}
\end{equation}
where $\boldsymbol{\Lambda} \coloneqq \operatorname{diag}(\lambda_1,
\lambda_2, \lambda_3)$
with eigenvalues $\lambda_1 \ge \lambda_2 \ge \lambda_3 > 0$, and
orthonormal columns
$\mathbf{v}_i \in \mathbb{R}^3$ define the principal axes of the
confidence ellipsoid.

\subsubsection*{Confidence Ellipsoid Radii}
At confidence level $1 - \alpha$ (e.g., $\alpha = \num{0.05}$ for a
$\qty{95}{\percent}$ confidence volume),
the semi-axis lengths $a_i$ of the spatial error ellipsoid are given by:
\begin{equation}
  a_i \coloneqq \sqrt{\lambda_i \chi^2_3(1 - \alpha)}, \quad i \in \{1, 2, 3\},
  \label{eq:supp_ellipsoid_radii}
\end{equation}
where $\chi^2_3(\num{0.95}) \approx \num{7.815}$ is the critical
quantile of the chi-squared distribution
with three degrees of freedom.

The horizontal confidence ellipse projection has semi-major axis
$a_h$ and semi-minor axis $b_h$:
\begin{equation}
  a_h \coloneqq \sqrt{\lambda_{h,1} \chi^2_2(1 - \alpha)}, \quad
  b_h \coloneqq \sqrt{\lambda_{h,2} \chi^2_2(1 - \alpha)},
  \label{eq:supp_horizontal_ellipse}
\end{equation}
where $\chi^2_2(\num{0.95}) \approx \num{5.991}$, and
$\lambda_{h,1}, \lambda_{h,2}$ are the eigenvalues
of the $2 \times 2$ horizontal sub-covariance matrix.

% =============================================================================
% 7. MAGNITUDE RESIDUAL PROPAGATION AND LOGARITHMIC ATTENUATION
% =============================================================================

\subsection*{Magnitude Residual Propagation and Logarithmic
Attenuation Calibration}
\label{subsec:supp_magnitude_residuals}

Local magnitude discrepancies $\Delta M_{\mathrm{L}, k} \coloneqq
M_{\mathrm{L}}(T_k) - M_{\mathrm{L}}(B_k)$
isolate differences between the automated synthetic Wood-Anderson
amplitude response
and the historical bulletin calibration.

\subsubsection*{Analytic Residual Expansion}
Expanding the Richter local magnitude definition for both solutions:
\begin{align}
  M_{\mathrm{L}}(B_k) &\coloneqq \log_{10} A_{\mathrm{WA}}(B_k) -
  \log_{10} A_0(r_{B,k}) + C_{S, B},
  \label{eq:supp_ml_bulletin} \\
  M_{\mathrm{L}}(T_k) &\coloneqq \log_{10} A_{\mathrm{WA}}(T_k) -
  \log_{10} A_0(r_{T,k}) + C_{S, T},
  \label{eq:supp_ml_target}
\end{align}
where $A_{\mathrm{WA}}$ is peak horizontal zero-to-peak amplitude in
millimeters on a standard
Wood-Anderson torsion seismometer ($T_0 = \qty{0.8}{\second}$, $h =
\num{0.8}$, $V_0 = \num{2080}$),
$r$ is hypocentral distance, $-\log_{10} A_0(r)$ is the empirical
regional attenuation function,
and $C_S$ is the station site correction.

Subtracting \cref{eq:supp_ml_bulletin} from \cref{eq:supp_ml_target}:
\begin{equation}
  \Delta M_{\mathrm{L}, k} \equiv
  \log_{10}\left(\frac{A_{\mathrm{WA}}(T_k)}{A_{\mathrm{WA}}(B_k)}\right)
  - \left[ \log_{10} A_0(r_{T,k}) - \log_{10} A_0(r_{B,k}) \right]
  + (C_{S, T} - C_{S, B}).
  \label{eq:supp_ml_residual_decomp}
\end{equation}
Under identical station selections ($C_{S, T} = C_{S, B}$) and
proximate hypocenters ($r_{T,k} \approx r_{B,k}$):
\begin{equation}
  \Delta M_{\mathrm{L}, k} \approx
  \log_{10}\left(\frac{A_{\mathrm{WA}}(T_k)}{A_{\mathrm{WA}}(B_k)}\right).
  \label{eq:supp_ml_approx}
\end{equation}
Consequently, a systematic magnitude bias $\mu_{\Delta
M_{\mathrm{L}}} \ne 0$ directly diagnoses
systematic amplitude scaling discrepancies in instrument
deconvolution or Wood-Anderson transfer function
implementation.
